\documentclass[output=paper]{langscibook}
\ChapterDOI{10.5281/zenodo.22796668}
\author{Tomáš Svoboda\affiliation{Charles University} and Olga Nádvorníková\affiliation{Charles University} and Věra Kloudová\affiliation{Charles University}}
\title{Introduction}
\abstract{\noabstract}
\IfFileExists{../localcommands.tex}{
  \addbibresource{../localbibliography.bib}
  \usepackage{tabularx,multicol}
\usepackage{url}
\urlstyle{same}

 
\usepackage{langsci-optional}
\usepackage{langsci-lgr}
\usepackage{langsci-gb4e}

\usepackage{paralist}
\usepackage{enumitem}
\usepackage{cprotect}
\usepackage{fancyvrb}

\let\clipbox\relax
\usepackage[export]{adjustbox}
\usepackage{array}
\usepackage[normalem]{ulem}
\usepackage{langsci-branding}
\usepackage{langsci-bidi}
\usepackage[linguistics]{forest}

  
\makeatletter
\let\thetitle\@title
\let\theauthor\@author
\makeatother

\newcommand{\togglepaper}[1][0]{
   \bibliography{../localbibliography}
   \papernote{\scriptsize\normalfont
     \theauthor.
     \titleTemp.
     To appear in:
     E. Di Tor \& Herr Rausgeberin (ed.).
     Booktitle in localcommands.tex.
     Berlin: Language Science Press. [preliminary page numbering]
   }
   \pagenumbering{roman}
   \setcounter{chapter}{#1}
   \addtocounter{chapter}{-1}
}

\newbool{bookcompile}
\booltrue{bookcompile}
\newcommand{\bookorchapter}[2]{\ifbool{bookcompile}{#1}{#2}}

\newcommand{\qual}[1]{`#1'}
\newcommand{\bv}[1]{(\ref{#1})}

\newcommand{\ip}{ImPost}
\newcommand{\ipadv}{ImPost-adverbial}
\newcommand{\pradv}{ProxPost-adverbial}
\newcommand{\chisq}{Chi-squared}


% for rotated columns
\newcolumntype{T}[2]{%
    >{\adjustbox{angle=#1,lap=\width-(#2)}\bgroup}%
    l%
    <{\egroup}%
}
\newcommand*\rot{\multicolumn{1}{T{45}{1em}}}%


\VerbatimFootnotes


% \newcommand{\keywords}[1]{\textbf{#1}}


%  \newcommand{\textarabic}[1]{\RL{\arabicfont #1}}
 
  %% hyphenation points for line breaks
%% Normally, automatic hyphenation in LaTeX is very good
%% If a word is mis-hyphenated, add it to this file
%%
%% add information to TeX file before \begin{document} with:
%% %% hyphenation points for line breaks
%% Normally, automatic hyphenation in LaTeX is very good
%% If a word is mis-hyphenated, add it to this file
%%
%% add information to TeX file before \begin{document} with:
%% %% hyphenation points for line breaks
%% Normally, automatic hyphenation in LaTeX is very good
%% If a word is mis-hyphenated, add it to this file
%%
%% add information to TeX file before \begin{document} with:
%% \include{localhyphenation}
\hyphenation{
    par-a-digm
}

\hyphenation{
    par-a-digm
}

\hyphenation{
    par-a-digm
}
 
  \togglepaper[1]%%chapternumber
}{}

\begin{document}
\maketitle 
%\shorttitlerunninghead{}%%use this for an abridged title in the page headers


Translation studies, founded in the 1970s, has become a thriving and autonomous research discipline ever since. This development was further fueled by the emergence of Descriptive Translation Studies in the 1990s \citep{Toury1995}, as well as the so-called turns that solidified the field’s merit \citep{Snell-hornby2006}. The significant evolution of parallel and comparable corpora in the 1990s (\citealt{Baker1993}, \citealt{Johansson1998}) allowed for enhancing empirical research in the field and the rise of corpus-based translation studies (CBTS). The current development of translation studies is evidenced by both the wealth of publications\footnote{{For an indicative overview, see the Translation Studies Bibliography –} \url{https://benjamins.com/online/etsb}{; the BITRA database} \url{https://aplicacionesua.cpd.ua.es/tra_int/usu/buscar.asp?idioma=en}{; or journals such as} {\textit{Meta, Perspectives,} }{and} {\textit{Target}}.} and a vivid research community that regularly convenes at well-established conferences. Among them, it is important to cite UCCTS{ (\textit{Using Corpora in Contrastive and Translation Studies)}, reflecting the close connection between CBTS and contrastive linguistics (see \citealt{BernardiniFerraresi2023}, \citealt{CzuloHansen-Schirra2017}), and especially the \textit{Translation in Transition} series, initiated in 2014 in Copenhagen, focusing on empirical approaches to translation studies. The idea for the present volume emerged from the sixth edition of this conference, held in Prague in  {September 2022}.}

The remarkable evolution of translation studies is largely propelled by its inherently interdisciplinary nature as rightly noted by \citealt{GrangerLefer2022}. Traditionally, the field has been associated with sub-fields and methodologies such as literary translation, the sociology of translation, translator studies, hermeneutics, legal translation studies and others.\footnote{For a first overview and details on the above see \citet{BakerSaldanha2019}.} For some time now, it has seen its applied fields (such as translator training, audiovisual translation, and machine translation) partially converge with empirical methodologies, including corpora and psycholinguistic experiments. Moreover, this interdisciplinary approach is often linked with broader linguistic frameworks, reflecting the socio-cultural and cognitive nature of translation, i.e. its product and its process, respectively. For our purposes here, among these frameworks the following are of special importance: cognitive linguistics and psycholinguistics (e.g., \citealt{Halverson2017}, \citealt{Carl2021}, and this volume), construction grammar (e.g., \citealt{CzuloEtAl2023}, etc.), systemic functional linguistics (e.g., \citealt{NeumannEtAl2017}), and variationist sociolinguistics (e.g., \citealt{PrieelsSutter2018}). The two conference series mentioned above exemplify this interdisciplinary potential, and the present volume seeks to highlight this aspect as well. The interdisciplinary approach was already highlighted in the preceding volume based on the \textit{Translation in Transition} conference held in Barcelona in 2019 (see \citealt{Bisiada2021}).

The present publication consists of ten chapters: this introduction by the editors and nine separate contributions. In total, seventeen authors contributed to the volume. {The core content originates from the 2022 \textit{Translation in Transition} conference mentioned above, showcasing the methodological and interdisciplinary diversity in contemporary translation studies. This is complemented by the opening and closing contributions that were not presented at the conference but address two current trends of translation studies. The opening chapter explores the research potential of advanced statistical methods and data visualization in the analysis of translation data in multilingual corpora, which are inherently multidimensional (see also   \citealt{DeSutterLefer2019}). The closing chapter addresses the growing field of interpreting studies, particularly in the context of the technological turn (}\citealt{Jimenez-Crespo2020}{), highlighting the paradigm shift brought about by ubiquitous machine translation, currently also generative AI, and enabled by evolving translation technologies.}

\section{General overview}

Across all chapters of the volume, interdisciplinarity is clearly a common thread. Readers are taken on a journey from a methodology focus (Chapter 2) through corpus-based translation studies (Chapters 4, 6–7), corpus-based contrastive linguistics (Chapter 5 and 7), audiovisual translation (Chapter 8) and specialized translation (nanomedicine in Chapter 6; legal translation in Chapter 7) to machine translation (Chapters 9 and 10). In Chapter 3, a cognitive approach to translation appears, adding a process-oriented viewpoint on translation to the pool.

Another unifying factor of all chapters is the coverage of novel topics, approaches, and methods. Each author brings a fresh perspective to their respective fields of study. At the same time, the novelty in each case is expressed differently. Some chapters, such as Chapter 2, introduce innovative methods in translation research; others offer new approaches in terms of thematic focus. Chapters 8 and 9 cover a novel theme in the study of audiovisual translation and machine translation, respectively. The overlapping interest in translation studies among the chapters goes hand in hand with a shared concern about how translation theory intersects with practical application (see particularly Chapters 3, 6, and 7, which focus on translator training).

From the point of view of methodology, chapters dealing specifically with contrastive analysis (Chapters 8–10) present qualitative assessments of translation and interpreting. Chapters 4, 5, 7, and 9 predominantly employ a quantitative approach to analyze linguistic data patterns and link those to translation research.

Beyond that, the volume covers a wide range of languages, including Arabic, Czech, English, French, German, Italian, Simplified Chinese, and Spanish – with a variety of language pairings and translation directions.

\section{Outlines of individual chapters}

The second chapter, written by Volker Gast and Vahram Atayan, deals with data visualization in translation studies as well as in corpus-based contrastive linguistics. Thus, it primarily represents a methodology-focused contribution. The authors begin by noting that the increasing availability of large multilingual corpora has created new opportunities for empirical research in these fields. They point out, however, that these “possibilities come with challenges” and that researchers need to filter relevant information and interpret complex data extracted from parallel (translation) corpora, in which the linguistic choices may be influenced by the input from the source language. The chapter aims to address these challenges by presenting and evaluating various visualization techniques.

After an overview of the importance of data visualization in understanding comparative linguistic data, the authors categorize comparative corpus data into three types: language-specific data, undirected comparative data, and directed comparative data. The chapter discusses three main visualization techniques: proportional visualizations, graph-based visualizations, and ordination techniques. Each technique is illustrated with concrete examples based on the authors’ previous research into so-called adverbials of immediate posteriority (IP-adverbials) in French and German. The authors also discuss the strengths and limitations of each technique and provide scripts (in R and Python) for readers to use in their own research.

One of the innovative aspects of the chapter is the introduction of “context-conditional translation graphs,” a new method for visualizing translation choices under specific contextual conditions. Arguably, this technique, along with Bayesian networks for modeling bilingual language production, represents a significant advancement in the field.

The third chapter, “Navigating translation effort through translation process research database to unravel dimensions of script and text direction” by Ali Saeedi, Longhui Zou, and Michael Carl, investigates the impact of script and text direction differences on the translation process, using eye-tracking and keylogging data to measure the technical, temporal, and cognitive effort exerted by translators. By analysing translations from English into Spanish, Simplified Chinese, and Arabic, the study reveals significant challenges posed by different scripts and text directions. The findings show that while script differences increase orthographic, lexical, and segment-level complexities, text-direction differences increase spatial orientation demands. Interestingly, no interaction effect between script and text-direction differences was observed, suggesting distinct cognitive processes for each challenge.

The study uses data from the CRITT Translation Process Research Database, focusing on from-scratch translation sessions. Indicators such as insertion, deletion, pause-to-word ratio, total number of fixations, and reading times are examined to gauge translation effort. The results indicate that translating between languages with similar scripts (English into Spanish) requires less cognitive effort than translating between languages with different scripts (English into Chinese and Arabic). Similarly, the temporal effort is higher for different-script translations, reflecting the additional cognitive load.

The findings of this study have significant implications for translation pedagogy, advocating for tailored training approaches that address script and text direction challenges separately. The study emphasizes the need for specific strategies to navigate these difficulties, such as exercises focusing on visual-spatial processing and spatial attention adaptation. By highlighting the nuanced cognitive demands of translation across different scripts and text directions, this chapter contributes valuable empirical insights and a deeper understanding of the translation process.

The fourth chapter, “A corpus-driven investigation of translation strategies for prenominal complex adjectives,” by Rossella Resi explores the translation of prenominal complex adjectives (PCAs) from German to Italian within a tourism-focused parallel corpus. Initially, German PCAs are categorized into seven groups: predicative structures with past and present participles, compound adjectives, multiple non-coordinated prenominal adjectives, adjectives modified by adverbs, adjectives modified by complex structures, and infinitive structures. The study uses regression analysis to assess the frequency of translation strategies, finding a strong correlation between PCA categories and translation strategies, with a notable tendency to preserve the original syntactic structures.

The chapter’s qualitative analysis focuses on the more transformational translation strategy of transposition, which involves changing grammatical categories without altering the meaning. Results indicate that the cognitive complexity inherent in German PCAs can encourage or trigger the use of transposition. Transposition, though less commonly applied, seems particularly effective for translating complex syntactic structures due to its ability to handle bidirectional syntactic reference and elevate semantic information within sentences beyond sentence level.

The findings suggest that while preserving original syntactic structures is a common strategy, transposition offers a viable alternative when dealing with intricate adjectival constructions. This highlights the importance of understanding the cognitive and syntactic challenges involved in translating PCAs and suggests that translation strategies should be adaptable to the specific complexities of the source text. Overall, the chapter provides valuable insights into the translation of complex adjectival structures and the application of different strategies to maintain meaning and coherence in the target language text.

Chapter five, “\textit{Schodiště z mramoru} and ‘the Marble Staircase’: English divergent counterparts of Czech postmodifying prepositional phrases,” is authored by Markéta Malá. This chapter focuses on noun-phrase modification, similar to the preceding chapter, but adopts a markedly different approach. In fact, large corpus-based studies have previously indicated that the majority of Czech postmodifying prepositional phrases, such as \textit{sochy z bílého mramoru} [statues of white marble], have a translation counterpart in English that employs a prepositional phrase as well (\citealt{KlégrMalá2009}). Using Johansson’s terminology \citep{Johansson2007}, widely adopted in contrastive linguistics, Malá thus asserts that most counterparts are \textit{congruent} rather than \textit{divergent}. This aligns with observations made about English, where postmodifying prepositional phrases represent the most common type of noun modification \citep{BiberEtAl2021}. Malá’s methodological innovation thus lies in her focus on the less frequent divergent counterparts.

The linguistic material used in this study is derived from a large multilingual corpus named InterCorp, which is limited to fiction (novels and short stories) and non-fiction in the Czech-English language pair. Adhering to the methodological requirements of translation studies, Malá meticulously considers the direction of translation in her analyses: she separately examines Czech postmodifying prepositional phrases in original Czech texts compared to their English translations (CZ-EN) and in Czech translations from English source texts (EN-CZ). A two-step analysis enables Malá to derive significant insights into the impact of register (text type) and direction of translation. However, the primary contribution of her study is the fine-grained formal and semantic typology of divergent English counterparts to the Czech postmodifying prepositional phrases. Her wealth of examples paves the way for future research in both contrastive linguistics and translation studies.

Chapter six, authored by Cécile Cuchet and Pascaline Dury, explores short diachronic corpora as a tool while translating emerging and new specialized fields, particularly in medical translation. It argues that diachronic corpora – texts collected over a few decades – can provide translators with valuable conceptual and terminological information to understand the evolution of terms in rapidly developing scientific fields, such as nanomedicine for cancer treatment. The authors performed a study using a corpus of scientific texts published in this domain from 1993 to 2017. The study aimed to identify useful diachronic information, such as term variations and the emergence of new terms, to aid translators in understanding complex, specialized terminology that is still evolving. The research highlights the challenges translators face when dealing with unstable terminology in emerging fields, where traditional lexicographic resources may prove insufficient.

In terms of methodology, the study emphasizes simplicity in corpus analysis to ensure that the approach is accessible to practicing translators, who may lack the time or technical expertise for complex analyses. The findings suggest that diachronic information can provide translators with a “minimal cognitive knowledge” necessary for translating new scientific concepts.

This approach holds potential for applications in translator training and other specialized fields where understanding the historic evolution of terminology is crucial. It also opens new avenues for research in translation studies, particularly in how translators can use diachronic data to improve translation accuracy in rapidly changing scientific domains.

Chapter seven, “Legal language in employment contracts and offer letters” by Patrizia Giampieri, investigates the differences in legal language and style between employment contracts and employment offers. Through constructing and analysing two specific corpora — one for contracts and the other for offer letters — the study aims to determine the extent to which these documents employ legalese or plain English. Legalese, characterized by complex prepositions, passive constructions, and polysemy, is compared to plain English, which strives for simplicity and accessibility. Findings suggest that employment offer letters use plain English more frequently, whereas employment contracts predominantly feature legalese, indicating different levels of formality and drafting styles.

The chapter covers the intricacy of legal language and the efforts by the Plain English movement to simplify the former. It highlights the use of corpora in legal language research, noting the challenges of accessing private legal documents and the necessity of DIY corpora for specific studies.

From the point of view of its methodology, the chapter describes the manual creation of corpora and the use of AntConc software for detailed linguistic analysis. Results show significant differences in legalese patterns between the two document types, such as the use of ‘shall,’ passive forms, and deictic elements.

There are considerable implications for translation studies, translator practice as well as training in that the research demonstrates how corpus-driven approaches can inform translation practices by highlighting language nuances and patterns. The chapter advocates for the application of rigorous corpus-based methods across disciplines to enhance understanding and translation of legal texts. Future research directions include exploring legal language in non-native contexts and analysing other types of legal documents.

The eighth chapter is entitled “Translating Rick \& Morty’s audiovisual novum parody: strategies and effects” and its author is Alice Ray. The chapter studies the popular genre of audio-visual translation, specifically the postmodern science fiction TV show originally appearing in English. The research analyses the series’ neologisms together with their counterparts in the French dubbed version. It considers the multi-faceted constraints of audiovisual translation, one of which is the relationship between text and image, focusing on parody and inherent humor in the combination of textual elements and visual content. The research applies a hybrid model involving contrastive linguistic methods and translation studies approaches.

The results of the analysis allow the reader to gain a better understanding of the challenges involved in translating what is referred to as ‘novum parody’ – the fictive terms that parody the sci-fi genre, resulting in humorous coinages. As a result of the complexity of the task, translators often have to prioritize some aspects of the audiovisual message contained in the invented terms over other aspects. Research findings involve the observation that translation of novum parody must be studied with a view to both narrative and visual contexts.

The analysis model put forward by the author is subsequently validated by way of analysing \textit{Futurama}, which is another audiovisual work under scrutiny here. The chapter is a valuable addition to the volume in terms of both methodology and field as it focuses on a niche topic, namely that of humor in sci-fi TV series. Apart from that, it presents a customized analytical model that can become useful in similar future research.

The final two chapters present research drawn from a similar triad of subjects, as they both compare the respective performance of professionals, trainees, and neural machine translation (NMT) systems. Whereas Kamenická’s chapter (nine) focuses on translation, Kloudová et al. (Chapter ten) focus on both interpreting and real-time MT of speech transcripts.

Renata Kamenická’s chapter, entitled “‘Default translation’ \textit{–} how much free of fault? An exploratory study of NMT and human translation comparison,” compares human and machine translation output. It looks into cognitive modeling of translation and, at the same time, assesses NMT performance in terms of how successful the latter is at imitating human translation. Crucially, the chapter employs and adapts Halverson’s concept of default translation as a relatively routinized and uninterrupted mode of translating – a concept touching upon the widely used notion of translation universals.

As for methodology, the performance of four NMT engines is compared to that of six translation trainees and six professional translators on texts in the English-Czech language pair. Translations are compared in both quantitative terms – using a custom metric – and qualitative terms as well.

From among the MT engines, DeepL was found to approximate human translation most successfully at the level of collocations. The chapter also puts forward a hypothesis that concerns the division of attention and effort of human translators. It stipulates that in their default mode – in order to avoid literalness – translators often favor translation solutions that are more readily available in terms of cognitive load, yet less literal than the ones considered to be optimal.

The main goal of the last chapter is well summed up in its title: “A comparison of human simultaneous interpreting and automatic simultaneous speech translation as a basis for further research.” The four authors, Věra Kloudová, Petra Mračková-Vavroušová, Dávid Javorský, and Ondřej Bojar, summarize the complexities of simultaneous interpreting (SI) before covering automatic simultaneous speech translation (ASST) and comparing these two services.

The study analyzes translation departures in SI on the one hand and ASST systems on the other. For SI, three interpreters were recruited; for the translation system, ELITR – a custom speech translation system – was used. Two speeches in English were analyzed as source texts, the target language being Czech in both cases. For the purpose of analysing content quality, a hybrid error-based quality assurance (QA) model was implemented.

The results show that the two trainee interpreters surprisingly outperformed a professional one. Compared to the interpreters the ASST model translated in a more literal fashion and its output involved fewer omissions and other translation departures than human interpreted renderings.

The challenges of the research involved the fact that two different services are compared here (interpreting activity results in a spoken output, whereas the output of ASST is a translated transcription of the original), and the focus on content in the comparison inevitably leads to partial results. The evaluation of omissions proved to be a challenge, too, knowing that in high-quality SI they are part of a deliberate strategy to produce a user-friendly rendering. The reflection of these research limitations represents a valuable methodological aspect of the chapter in its own right. With its well-resourced summary of evaluating interpreting quality, its methodological focus, and its novel thematic coverage, the chapter complements the other approaches well in that it provides a foundation for further research.

\section{Research limitations and potential}

The research limitations of this volume logically stem from its greatest asset – interdisciplinarity. Tackling a great many areas often sacrifices depth for breadth. We believe that we have struck the right balance between both. Another challenge of the volume’s agenda has been the dynamic nature of some areas covered, such as emerging methodologies and the development of machine translation engines. Furthermore, the chapters in this volume were written largely before the rapid emergence of generative artificial intelligence and therefore do not address explicitly this topic, which has attracted much attention lately. Nevertheless, regardless of future technological developments, the interdisciplinary nature of translation studies, reflected in the contributions to this volume, is bound to remain one of the field’s greatest strengths and a key foundation for its continued relevance and development.

 Beyond that, the vitality of the field is evident in the diversity of topics covered in the forthcoming \textit{International Encyclopedia of Language and Linguistics} (Nesi \& Milin, forthcoming), which spans 78 different areas, across domains such as localization, translation policy, collaborative translation, cognitive aspects in translation, and neurolinguistics in translation studies. New perspectives are also emerging through innovative types of corpora, such as the MUST corpus (\citealt{GrangerLefer2020}), which offers a wealth of metadata for translation studies research, and multilingual translation corpora enriched with specific linguistic information – for example, the InterCorp multilingual corpus includes Universal Dependencies annotation and syntactic complexity measures \citep{RosenEtAl2024}. However, as highlighted during the round table discussion at the Translation in Transition conference in Prague in 2022, the future of translation studies will largely depend on researchers’ ability to comprehend the theoretical grounding of the domain, allowing them a better understanding of the nature of translation and the sociocultural, technological, and cognitive factors that shape it. We hope to have added yet another example of the interdisciplinary potential of translation studies and provided stimuli in search of its common ground in the future. It may well be that this common ground will involve a rediscovery of some grand and ambitious attempts at charting the territory by some of the founding figures of the discipline, such as Jiří Levý.  {In 1965}, in respect to literary theory, making a subsequent link to translation studies, he had the following to say:

\begin{quote}
Currently, particularly favorable conditions have emerged for the refinement of literary scholarship, since most of the adjacent social sciences have come within the reach of \textit{exact methods}. … The sciences that achieve far more exact and indisputable results [compared to speculative statements of some literary historians] formulate their unembellished conclusions from a specific set of observations, and only upon those conclusions they formulate further hypotheses, thus creating a \textit{chain of increasingly adequate hypotheses}. \citep[11, 18, our translation; emphasis added]{Levy1965}
\end{quote}

Similarly to the early efforts undertaken at the dawn of translation studies in mid-1960s, today as well, interdisciplinarity and respect for exact methods still appear to be two of the most promising directions for our interdiscipline.

\section*{Acknowledgements}

The editors of this volume wish to thank all co-authors involved for their contributions and patience in the process of producing the book, to the reviewers for their input as well as Language Science Press representatives and editors for their support and expert advice.
The production of this publication was supported by \textit{Cooperatio lingvistika} project of Charles University’s Faculty of Arts in 2024.

\sloppy\printbibliography[heading=subbibliography,notkeyword=this]
\end{document}
